\documentclass{article}
\usepackage{amsmath}
\usepackage{amsfonts}
\usepackage{amssymb}
\usepackage{graphicx}
\usepackage{hyperref}
\usepackage[a4paper, margin=0.5in]{geometry}

\title{Lander Challenge Summary}
\author{Yoonjae Hwang}

\begin{document}

\maketitle
\section{Introduction}
In this project I explored geometric structure preserving numerical integration methods, 
and geometric control systems. 
Lie Group Variational Integrators finds discrete equations of motion
by discretising the Lagrangian of the system, and applys the discrete case for the variation
of the system needing to be 0. 
The spacecraft position and its attitude is represented by the configuration manifold given by 
$\mathbf{SE}(3)$.
\subsection{Motivation}
\begin{itemize}
    \item $SE(3)$ and $SO(3)$ are Lie Groups. 
    \item Certain quantities are conserved
    \item Runge Kutta Schemes do not exhibit these quantities preserving, thus resulting in poor long time accuracy in qualitattive behaviour
    \item By finding discrete equations of motion that is computed directly on the Lie Groups, these quantities can be conserved.
\end{itemize}
These are Lie Groups.

\section{Equations of Motion}
For continuous

\end{document}